%%%%%%%%%%%%%%%%%%%%%%%%%%%%%%%%%
%
%       EXERCÍCIO
%
%%%%%%%%%%%%%%%%%%%%%%%%%%%%%%%%%

\definevalorquestao

\begin{exercicioBanco}[\valorquestao]
Resolva, em \(\bbR\), a inequação
%
\[
    \dfrac{4x^{2}+8x+13}{x^{2}-4x+3}\leq 1.
\]
%

% Define as alternativas
\newcommand{\alternativas}{%
    \begin{center}
        \begin{tabularx}{\textwidth}{XXX}
            (a) \(]-\infty,-1] \cup [3,\infty[\). &
            (b) \([-1,1[\ \cup\ ]3,\infty[\). &
            (c) \(]-\infty,-1] \cup ]1,3[\). \\[5pt]
            (d) \([-1,1[ \cup ]1,3[\). &
            (e) \([-1,3]\).
        \end{tabularx}
    \end{center}
}

% Define a resposta correta
\newcommand{\resposta}{D}

% Lógica condicional para exibição
\ifdefstring{\atividade}{lista}{%
    \alternativas
    \vspace{0.5em}
    
    \noindent\textbf{Resposta:} letra \textbf{\resposta}.
}{%
    \ifdefstring{\modo}{objetivo}{%
        \alternativas
    }{%
        % modo = discursiva → não mostra alternativas
    }
}
\end{exercicioBanco}

